\chapter{Jak rozwijac projekt od malego do sredniego serwisu}

\section{Etap 1: jedna strona i jeden chunk}
Kazdy projekt warto zaczac od minimalnego pionowego przekroju. Jedna strona, jeden chunk,
jeden punkt \mc{moon(...)}. Celem nie jest od razu architektura idealna, tylko pierwszy
poprawny wynik HTML.

\section{Etap 2: wspolne metadata}
Gdy pojawia sie druga strona, zwykle od razu warto wydzielic chunk z metadata wspolnymi.
To pierwszy realny krok w kierunku reuzywalnosci.

\section{Etap 3: funkcje pomocnicze}
Gdy te same obliczenia zaczynaja powtarzac sie w kilku miejscach, wydziel funkcje.
Na tym etapie MoonChunk zaczyna przynosic realny zysk: projekt nie tylko generuje HTML,
ale robi to w sposób uporzadkowany.

\section{Etap 4: dane z JSON}
Kiedy tresc strony przestaje byc recznie wpisywana, przychodzi pora na \mc{data(...)}.
To moment, w ktorym projekt przestaje byc tylko zestawem statycznych napisow, a zaczyna
stac sie malym systemem generowania dokumentow na podstawie danych.

\section{Etap 5: podzial na pliki}
Jesli plik glowny zaczyna byc zbyt duzy, rozdziel odpowiedzialnosci:

\begin{itemize}
  \item metadata,
  \item strony glowne,
  \item logika bloga,
  \item sekcje wspolne,
  \item dane.
\end{itemize}

W tym momencie importy staja sie naturalna częscią codziennej pracy.

\section{Etap 6: wzorce projektowe}
Kiedy projekt ma juz kilka stron i kilka rodzajow danych, warto swiadomie wprowadzic
powtarzalne wzorce:

\begin{itemize}
  \item sekcje budowane z tablic obiektow,
  \item funkcje etykietujace,
  \item wspolne metadata z lokalnym nadpisaniem,
  \item jawne punkty wejscia \mc{moon(...)} dla wybranych przebiegow generacji.
\end{itemize}

\section{Etap 7: czyszczenie i stabilizacja}
Na tym etapie warto przejrzec projekt i zadac kilka pytań kontrolnych:

\begin{itemize}
  \item czy nazwy sa nadal czytelne,
  \item czy logika nie rozlala sie do HTML-a,
  \item czy globali nie jest za duzo,
  \item czy importy oddaja prawdziwy podzial odpowiedzialnosci,
  \item czy debugowe \mc{print(...)} nadal sa potrzebne.
\end{itemize}

\section{Etap 8: swiadome skalowanie}
MoonChunk nie musi rosnac przez doklejanie kolejnych przypadkowych plikow. Dobrze prowadzony
projekt rozwija sie warstwowo. Każda nowa strona albo sekcja znajduje swoje miejsce w
istniejacej strukturze.

\section{Zasada, ktora ratuje projekt}
Najwazniejsza rada na etapie wzrostu jest prosta: \important{najpierw nazywaj i porzadkuj,
potem dopiero rozbudowuj}. W wielu projektach to wlasnie brak porzadku, a nie brak funkcji,
staje sie prawdziwym problemem.
